\newpage
\chapter*{\sffamily Resumen}
\addcontentsline{toc}{chapter}{Resumen}%

\par Las interfaces cerebro-computador (BCI) basadas en electroencefalografía (EEG) constituyen herramientas de gran potencial para la decodificación neuronal en aplicaciones clínicas, asistivas y cognitivas debido a su bajo costo, portabilidad y alta resolución temporal. Sin embargo, las señales EEG presentan baja especificidad espacial, alta susceptibilidad al ruido y los artefactos, dinámicas no lineales y no estacionarias, y una elevada variabilidad entre sujetos y dentro de un mismo sujeto. Estas características dificultan la robustez, la generalización y la interpretabilidad de los modelos actuales de aprendizaje profundo.

Esta investigación doctoral desarrolla un marco de aprendizaje profundo guiado por conectividad para la decodificación de señales EEG, mediante tres contribuciones metodológicas complementarias. La primera, EEG-GCIRNet, introduce una representación de conectividad funcional gaussiana para capturar dependencias espacio-espectrales entre canales. Las representaciones obtenidas se transforman en imágenes topográficas de EEG y se procesan mediante un autoencoder variacional que realiza conjuntamente reconstrucción y clasificación. En tareas de imaginación motora, EEG-GCIRNet alcanzó una exactitud promedio del 81,82\,\% en el escenario binario del conjunto GIGA MI y del 75,20\,\% en el escenario de cinco clases de PhysioNet EEGMMI. Los análisis de reconstrucción, espacio latente y mapas de relevancia permitieron examinar la organización espacial y discriminativa de las representaciones aprendidas.

La segunda contribución, \gls{tektenet}, integra filtrado temporal no lineal, reconstrucción del espacio de estados de Takens y estimación de entropía de transferencia basada en núcleos dentro de una arquitectura diferenciable. Este enfoque permite modelar dependencias predictivas no lineales, dirigidas y temporalmente retardadas entre canales EEG. Los experimentos con datos semisintéticos permitieron analizar su comportamiento frente a estructuras de interacción controladas y diferentes niveles de ruido, mientras que la evaluación sobre los datos oficiales de prueba de BCI Competition IV-2a alcanzó una exactitud del 80,12\,\%. Los análisis temporales, espectrales y de conectividad permitieron caracterizar los intervalos de tarea, las respuestas frecuenciales aprendidas y los patrones dirigidos asociados con la decodificación de imaginación motora.

La tercera contribución, CTE-Net, combina contextualización basada en Transformers, filtrado temporal no lineal, reconstrucción de Takens y estimación diferenciable de entropía de transferencia para aprender representaciones de conectividad dirigida en EEG pediátrico. El modelo fue evaluado para la clasificación de niños con trastorno por déficit de atención e hiperactividad (TDAH) y controles mediante particiones por sujeto y entrenamiento repetido. CTE-Net alcanzó una exactitud del 83,4\,\% y un área bajo la curva ROC del 90,2\,\% a nivel de participante. Los análisis de ablación respaldaron principalmente la contribución del Transformer y mostraron evidencia más moderada sobre el beneficio predictivo adicional de la entropía de transferencia. Asimismo, la representación de conectividad dirigida redujo la dispersión intraindividual en el espacio de análisis y permitió examinar la organización por clase y la estabilidad de las conexiones aprendidas.

En conjunto, los resultados muestran que las representaciones de conectividad funcional y dirigida pueden integrarse con modelos de aprendizaje profundo para combinar desempeño predictivo con el análisis explícito de relaciones espaciales, temporales y multicanal. No obstante, las conexiones estimadas a nivel de sensor se interpretan como dependencias estadísticas y predictivas dirigidas en el sentido de Wiener--Granger, y no como evidencia de causalidad mecanicista o anatómica. Asimismo, las tres contribuciones fueron evaluadas en diferentes conjuntos de datos, tareas y protocolos experimentales, por lo que sus resultados no constituyen una comparación controlada entre arquitecturas sucesivas.

Esta investigación se desarrolla en el marco del proyecto Alianza Científica con Enfoque Comunitario para Mitigar Brechas de Atención y Manejo de Trastornos Mentales Relacionados con Impulsividad en Colombia (ACEMATE) y contribuye al desarrollo de metodologías de EEG interpretables y conscientes de la conectividad. Su alcance se relaciona con el Objetivo de Desarrollo Sostenible 3, orientado a la salud y el bienestar, y el Objetivo de Desarrollo Sostenible 4, relacionado con una educación inclusiva y equitativa.

\vspace{2cm}
\textbf{Palabras clave}: electroencefalografía; interfaces cerebro-computador; aprendizaje profundo; conectividad funcional; conectividad efectiva; entropía de transferencia; imaginación motora; trastorno por déficit de atención e hiperactividad.